%% ma: L12-B04-0012
%% lop: 12
%% chuong: 1
%% bai: 4
%% dang: 12.4.5
%% loai: TN4
%% muc: TH
%% dap_an: "D"
%% nguon: "(NBV)-ĐỀ SỐ 15-MỖI NGÀY 1 ĐỀ THI 2025.pdf (D15, MỖI NGÀY 1 ĐỀ - Đề số 15), Phần I Câu 4"
%% kien_thuc: "Bài 3 (tiệm cận), Bài 4 (nhận dạng đồ thị hàm phân thức bậc nhất trên bậc nhất)"
%% trang_thai: chinh-thuc
%% ghi_chu: "Trùng ý tưởng (không trùng hoàn toàn) với L12-B04-0005, nhập vào cùng dạng 12.4.5. Đáp án gốc D đúng. Hình vẽ lại bằng y = (x+2)/(x+1) theo tiệm cận x = -1, y = 1 và giao điểm với hai trục."
\begin{ex}
Đường cong ở hình vẽ là đồ thị của hàm số nào trong các hàm số sau?
\begin{center}
\begin{tikzpicture}[scale=.75]
  \draw[phu] (-1,-3.4)--(-1,3.6);
  \draw[phu] (-4.3,1)--(4.3,1);
  \draw[truc] (-4.4,0)--(4.5,0) node[below]{$x$};
  \draw[truc] (0,-3.5)--(0,3.8) node[left]{$y$};
  \draw[thick,domain=-.62:4.2,samples=60,smooth] plot (\x,{(\x+2)/(\x+1)});
  \draw[thick,domain=-4.2:-1.45,samples=60,smooth] plot (\x,{(\x+2)/(\x+1)});
  \node[below right,font=\small] at (0,0) {$O$};
  \node[above left,font=\small] at (0,1) {$1$};
  \node[below,font=\small] at (-1,0) {$-1$};\end{tikzpicture}
\end{center}
\choice{$y=\dfrac{x-2}{x+1}$}{$y=\dfrac{x+2}{x-1}$}{$y=\dfrac{x-2}{x-1}$}{\True $y=\dfrac{x+2}{x+1}$}
\loigiai{Đồ thị có tiệm cận đứng $x=-1$, tiệm cận ngang $y=1$ nên loại B, C. Đồ thị cắt trục tung tại điểm có tung độ lớn hơn $1$ (cụ thể $y(0)=2$) nên loại A, chọn D.}
\end{ex}
